\section{Algebra}
\subsection{Kvadratsætningerne}
\begin{align*}
    (a+b)^2 &= a^2 + 2ab + b^2 \\
    (a-b)^2 &= a^2 - 2ab + b^2 \\
    (a-b)(a+b) &= a^2 - b^2
\end{align*}

\subsection{Andengradsligninger}

\subsubsection{Diskriminant af andengradsligning}
\[
    D = b^2 - 4ac
\]

\subsubsection{Rødder af andengradsligning}
En andengradsligning på formen \( a x^2 + b x + c = 0 \) har rødderne
\[
    x = \frac{-b \pm \sqrt{b^2 - 4ac}}{2a}
\]

Særtilfælde $c=0$
\[
    x = 0 \lor x = -\frac{b}{a}
\]

\subsubsection{Vendepunkt (vertex)}
En andengradsligning på formen \( y = a x^2 + b x + c \) har vendepunktet
\[
    \left( -\frac{b}{2a}; c - \frac{b^2}{4a} \right)
\]

\subsubsection{Retning af andengradsligning}
Hvis \( a > 0 \) vender parablen opad, og hvis \( a < 0 \) vender parablen nedad.

\subsubsection{Værdimængden af en funktion}
For en andengradsligning på formen \( y = a x^2 + b x + c \):
\begin{itemize}
    \item Hvis \( a > 0 \): \( [y_{\text{min}}, \infty[ \)
    \item Hvis \( a < 0 \): \( ]-\infty, y_{\text{max}}] \)
\end{itemize}
hvor \( y_{\text{min}} \) og \( y_{\text{max}} \) er y-værdien af vendepunktet.